\section{Conclusion}

This report introduces Spartan Write, a document authoring software powered by agentic AI.
It discusses the shortcomings of authoring documents with general-purpose word processing software and highlights the steep learning curve that \LaTeX-based solutions like Overleaf offer.

\subsection{Key findings and interpretation}
The benchmark turns up a somewhat surprising insight: differences in cost and latency among modern LLMs come down less to the raw skill or speed of the models, and more to the way they use tools.

All capable models complete tasks at roughly the same rate (Table~\ref{tab:task-completion}), but they get there using very different strategies.
For example, DeepSeek repeatedly reads files it doesn't need, which pushes its token use up to $62{,}485$ per run—$3.6\times$ what GPT-5.4 uses (Table~\ref{tab:token-averages})—and results in a mean duration of $148.6$\,s (Figure~\ref{fig:bench-duration-vs-tokens}).
This isn’t a problem with DeepSeek’s reasoning. Instead, it’s a side effect of tool-invocation habits that re-insert file contents into context on every turn.
Gemini~3~Flash, in contrast, uses a different approach: it restructures in-place using \texttt{rename\_file} (Figure~\ref{fig:bench-tool-use}) rather than reading and then editing, which delivers the lowest latency ($12.9$\,s) and the second-lowest cost ($\$0.0094$ per run, Figure~\ref{fig:bench-cost-per-model}).

This finding points to two main takeaways.
First, for document-editing workloads, models are best ranked by how well their tool use matches efficient patterns.
Second, prompt engineering and tool-schema design can nudge a model from being an expensive outlier to a competitive choice.
It’s notable that the GPT family never uses \texttt{rename\_file} in any of the $102$ completed runs (Figure~\ref{fig:bench-tool-use}), even when rename-based solutions are available. This suggests a schema-level incompatibility, rather than any lack of reasoning.

Across all models, the hardest tasks involve large-document structural changes (Table~\ref{tab:hardest-tasks}, Figure~\ref{fig:bench-duration-by-job}), not weaknesses tied to any specific model.
This is a signal that the true complexity of document editing comes from the nature of the work itself, not from the AI models. Improvements should focus on making large-context reasoning more general-purpose, rather than fine-tuning each model separately.

Based on these findings, \texttt{google/gemini-3-flash-preview} stands out as the recommended primary model for Spartan Write. It’s the fastest ($12.9$\,s), nearly the cheapest ($\$0.0094$ per run), and it aligns well with document-editing workflows.
That said, this recommendation is based on efficiency—not correctness. Completing a task quickly and cheaply doesn't always mean the result is the best possible.

\subsection{Contributions}

This work makes two contributions to the LLM-driven document authoring space.

First, the report presents Spartan Write, a document authoring system demonstrating how agentic AI reduces cognitive load in \LaTeX-based technical writing.
Rather than requiring users to master tool-specific syntax for editing, the system abstracts file operations and document structure behind a natural-language interface, letting users express intent and delegate implementation details to the agent.

Second, we contribute a benchmark methodology for evaluating LLMs on document-editing tasks.
The suite comprises 17 parameterised test cases spanning single-edit operations, large-scale refactoring, error correction, and multimodal tasks.
It includes rigorous metrics for reliability, latency, token usage, cost, and tool invocation patterns.

\subsection{Limitations and future work}

One key limitation is important to keep in mind: this benchmark measures \textbf{efficiency} rather than output \textbf{correctness}.
The bench reports task completion rates and other performance metrics, but “completion” here only means the agent ran without a runtime error.
It doesn't guarantee that edits will always match what the user wanted, or that the document will remain correct.
So, the recommendation for Gemini is based solely on efficiency, and it assumes that all models deliver equally correct results, which is an assumption.

This matters because if output quality varies significantly across models, the efficiency-based rankings could be misleading.
In practice, a model that's a bit slower but more reliable might be a better choice, even if the benchmark suggests otherwise.
A more complete evaluation of output quality would help fill this gap by specifying specific operations on a case-by-case basis.
For example, the content generation test could include an assertion that the write tool is never called.
Other possible approaches include comparing edits to reference versions, structured annotation, or asking an LLM to judge a sample of benchmark runs.
Adding these quality metrics is the main priority for future work, and would give more useful context for choosing a model.

Beyond quality assessment, two secondary research directions remain: investigating the GPT family's systematic avoidance of \texttt{rename\_file} and validating the benchmark against real-world usage through human-centred studies measuring task completion time, error rates, and user satisfaction with the agent-assisted workflow.

Several Spartan Write UX feature enhancements merit exploration: Undo support for agent actions, real-time multi-user collaboration, and integration of SyncTeX, which would allow the agent to edit documents more accurately by mapping PDF coordinates back to source locations.
